\documentclass[../../../main/main.tex]{subfiles}

\begin{document}

\chapter{Midpoints and Bisectors}

The \textit{midpoint} of a segment $AB$ is the point $M$ such that $AM \cong MB$ and $[AMB]$.

\begin{axiom}[Midpoint Existence]
    $\forall A,B\in\mathbf{P}, A\neq B \implies \exists M\in\mathbf{P}\colon (\mathcal{B}(A,M,B)) \land (AM \cong MB)$.
\end{axiom}

\begin{axiom}[Midpoint Uniqueness]
    $\forall A,B\in\mathbf{P}, A\neq B \implies \forall M,N\in\mathbf{P}\colon (AM \cong MB) \land (\mathcal{B}(A,M,B)) \land (AN \cong NB) \land (\mathcal{B}(A,N,B)) \implies M = N.$
\end{axiom}

https://chat.deepseek.com/a/chat/s/c1efa70e-bf74-4e92-b967-4054f3472f26

By removing betweenness as a criteria, we obtain the set of points $\{ X \mid AM \cong MB \}$. This set is a line that is perpendicular to the segment $AB$. Alternatively, we can define the perpendicular bisector as the line containing the midpoint $M$ of segment $AB$


\end{document}
